% R15 component; only input paths and enumerated location/grammar prose are edited.
% Source: revisions/2026-10-04-r15/manuscript/implementation_summary.tex; commit 1cb3cc9efe135966ec228dfaf51c6841a6dfc97c.
\subsection{Information and the implemented decision rule}\label{sec:r15implementation}
The economic institution is a controller that receives capital observations at specified dates, stores an internal state, and holds the selected consumption share until the next decision date. Its economic primitive is the original capital diffusion. Its internal numerical state need not equal the physical state. When the controller uses private randomization, that randomization is part of its specified information and is included in the law of the evaluated policy.

The exact-history representation in Proposition~\ref{prop:r12observation} assumes continuous noiseless observation of every capital coordinate, known coefficients and actions, exact drift integration, and an independent private Brownian motion. Under those assumptions the observation history supplies a Brownian representation with the law required by the innovation-driven controller. It does not recover the original latent null-space shock path by path. Proposition~\ref{prop:r14randomization} states why the added independent randomization leaves the economic optimum unchanged. These are representation results for a randomized history-dependent policy.

Finite observations require a different implementation account. Proposition~\ref{prop:r14sensor} bounds the difference between the measured-state controller and its specified protected parent, retaining sensor error, the internal recurrence, action-map error, clipping, and arithmetic. Its allowance belongs to that parent and that time grid. The empirical comparison reports it alongside the gain measured for the same parent and population; a certificate for a different mesh is not substituted. Complete assumptions and proofs are in Appendix~\ref{app:r15observationtheory}.

